\chapter{维数}\label{chap:dimension}

\noindent\textbf{第 \thechapter{} 章学习目标}
\begin{itemize}
  \item 从直观上描述仿射簇的维数, 并用坐标环中的代数独立性描述它.
  \item 定义仿射簇的函数域, 并在例子中求出函数域.
  \item 用超越基度量域扩张.
  \item 在例子中仔细算出仿射簇的维数.
  \item 用 Noether 正规化定理证明维数论基本定理.
\end{itemize}

给定一组多项式方程, 关于它的解集, 最基本的几何问题大概是这个解集有多\textquotedbl{}大\textquotedbl{}. 本章引入仿射簇维数这个重要概念, 以说明这个问题.

直观上, 维数衡量我们在一个集合中移动的自由度. 那么, 怎样度量在仿射簇中移动的自由度? 我们将看到, 回答这个问题的关键在坐标环里. 本章的一个出发点是: 仿射簇的维数就是其坐标环中代数独立元素的最大个数.

这也许会让读者想起线性代数中的维数: 一个 (有限维) 向量空间的维数可以定义为线性独立元素的最大个数. 本章所发展的许多思想, 确实都与线性代数中关于维数的思想类似. 尤其是第~\ref{sec:transcendence-bases}~节引入的超越基, 其定义和性质与向量空间的基十分相近. 为了建立完善的超越基理论, 我们需要在域而非环中工作. 因此第~\ref{sec:function-fields}~节专门讨论不可约仿射簇的函数域.

在建立函数域和超越基的基础之后, 第~\ref{sec:dimension-of-affine-varieties}~节将定义维数, 并描述它的一些基本性质. 例如, 我们会看到 $\dim(\mathbb{A}^{n})=n$; 若 $X\subseteq Y$, 则 $\dim(X)\leq\dim(Y)$. 本章的思想最终汇聚于第~\ref{sec:fundamental-theorem-of-dimension-theory}~节, 得到我们所称的维数论基本定理. 这个代数几何中的核心结果大意是: 一组多项式方程中, 每增加一个方程, 解集的维数至多降低一.

\par\smallskip
{\footnotesize
\noindent\textcopyright{} The Author(s) 2025\par
\noindent E. Clader, D. Ross, \emph{Beginning in Algebraic Geometry}, Undergraduate Texts in Mathematics,
\href{https://doi.org/10.1007/978-3-031-88819-9_7}{\nolinkurl{https://doi.org/10.1007/978-3-031-88819-9_7}}.\par}
